
\chapter{Introduction}

\section{Background}

Containers have become the dominant deployment model for modern applications. Unlike virtual machines, which provide hardware-enforced isolation, containers share the host kernel with every other container on the node, which concentrates a substantial attack surface in a single component~\cite{wang2023}. A single kernel vulnerability can therefore compromise every container co-resident on the host~\cite{cve2019}. Runtime-engine flaws add a distinct escape path: CVE-2019-5736 overwrites the host \texttt{runc} binary via writable \texttt{/proc/self/exe} during container attach~\cite{cve2019,unit42cve20195736}, and Reeves et al.~\cite{reeves2021runtime} show that many other published container escape proofs-of-concept against engines including \texttt{runc} stem from host interfaces leaked into the container namespace, not only from kernel bugs.

The Open Container Initiative (OCI) standardises container runtimes, and its reference implementation, \texttt{runc}, is intentionally biased toward performance and compatibility rather than confinement~\cite{volpert2024}. Linux already provides the kernel primitives that would close the gap: seccomp filters at the syscall boundary~\cite{lopes2020}, mandatory access control through AppArmor, and Linux capabilities to split the historical \texttt{root} into smaller privileges. In practice, however, all three are only as strong as the policies operators ship with their workloads, and meaningful policies are tedious and error-prone to author by hand.

The industry has reacted in three directions. Sandboxed runtimes such as gVisor and Kata Containers replace the shared-kernel model with either a user-space kernel or a per-container microVM, trading 20--100\% runtime overhead for stronger isolation~\cite{wang2022,viktorsson2020}. Cluster-side projects such as the Kubernetes Security Profiles Operator (SPO)~\cite{spo2025}, Inspektor Gadget~\cite{inspektorgadget2025}, and Tetragon~\cite{tetragon2025} record container behaviour through eBPF and produce seccomp/AppArmor profiles that can later be pinned to pods. Both directions raise the floor of container security significantly, but both also introduce non-trivial cost: a sandbox costs CPU, while a recorder costs cluster infrastructure plus an engineer who knows how to operate it.

\section{Problem Statement and Research Gap}

This thesis addresses three concrete gaps that emerge from the literature and from the state of the art surveyed in Chapter~\ref{ch:litreview}.

\textbf{Gap 1: hardened \texttt{runc} is under-evaluated.} Prior work measures default \texttt{runc} configurations and reports weak isolation relative to sandboxes, but does not evaluate what a properly hardened shared-kernel runtime achieves when seccomp, AppArmor, and capabilities are applied together from workload-specific profiles. Existing runtime comparisons~\cite{wang2022,volpert2024,viktorsson2020} therefore leave open whether profile-driven confinement on \texttt{runc} can materially narrow the gap to sandboxed alternatives without replacing the runtime stack.

\textbf{Gap 2: no confinement tooling without orchestration.} Operators who run containers with plain \texttt{runc}---without Kubernetes or a comparable control plane---must either hand-write seccomp, AppArmor, and capability policies (hours per service, deep kernel expertise, brittle across dependency updates) or accept the stock defaults. Cluster-side recorders~\cite{spo2025,inspektorgadget2025} and mechanism-specific research prototypes~\cite{lopes2020,canella2021} do not deliver an integrated, runtime-local workflow invokable from a developer workstation.

\textbf{Gap 3: limited integration across mechanisms.} Most prior automated approaches~\cite{lopes2020,canella2021} target seccomp alone. AppArmor, capability narrowing, and the orchestration of all three under one workflow remain fragmented across separate utilities and bundles.

\section{Research Questions}

\begin{enumerate}
    \item Can the runtime itself, without infrastructure changes, generate seccomp, AppArmor, and capability profiles whose security level approaches what a skilled operator would produce by manual tuning or what a cluster-side recorder produces in a CI pipeline?
    \item How should a runtime expose this capability so that it remains a one-flag operation for a developer, while still satisfying the constraints that make dynamic profiling trustworthy (child-process coverage, profile relaxation during the recording, non-root recommendation)?
    \item After scan-then-enforce on a representative workload, how much does post-exploitation containment improve relative to stock \texttt{runc} with default policies?
\end{enumerate}

\section{Proposed Approach}

The proposed runtime is a fork of \texttt{runc} that adds a single flag, \texttt{-{}-security-scan}. When the flag is present:

\begin{itemize}
    \item All host-supplied confinement is relaxed in memory for the run (seccomp cleared, \texttt{NoNewPrivileges} disabled, all known capabilities granted, AppArmor swapped for a complain-mode profile). The on-disk \texttt{config.json} is not mutated during the run.
    \item Syscalls are recorded by \texttt{oci-seccomp-bpf-hook} and emitted as \texttt{generated/seccomp.json}.
    \item Capabilities are traced cgroup-wide by \texttt{capable-bpfcc} with a \texttt{bpftool}-pinned cgroup map, so child processes are observed.
    \item AppArmor learns through a complain-mode profile loaded by a \texttt{createRuntime} hook and unloaded by a \texttt{poststop} hook.
    \item All hooks are sub-commands of the runtime binary itself (\texttt{runc scan-aa-load}, \texttt{runc scan-cap-trace-start}, etc.). The OCI bundle never has to embed executable artefacts.
    \item On finalisation, the runtime narrows --- never widens --- \texttt{process.capabilities} in \texttt{config.json} to the observed set.
\end{itemize}

A second mode, \emph{enforcement}, runs the workload with the generated profiles applied. The implementation also revisits the default capability set and ships two opt-in flags (\texttt{-{}-default-capabilities} for the upstream three-cap set, \texttt{-{}-default-capabilities-docker} for the historical fourteen-cap set), so that the hardened default does not break existing bundles.

\section{Experimental Evaluation}

Evaluation has two parts. \textbf{Empirical.} The \texttt{runc-hardened-test} laboratory exercises the scanner on an intentionally vulnerable Flask workload: it validates that scan-then-enforce produces usable profiles, documents edge cases and drawbacks (file capabilities, complain-mode audit parsing, scan coverage), and compares post-exploitation containment against stock \texttt{runc} on the same injection payloads. \textbf{Analytical.} Chapter~\ref{ch:experiments} also situates the runtime against manual authoring, cluster-side recorders, and sandboxed runtimes using the competitive landscape from Chapter~\ref{ch:litreview} (Table~\ref{tab:competitors}) and published measurements from prior work---without running sandbox benchmarks on the fork itself. Quantitative performance comparison against gVisor, Kata, or a dedicated benchmark harness is deferred to future work.

\section{Expected Contribution}

\begin{enumerate}
    \item A practical, developer-grade runtime that produces seccomp, AppArmor, and capability profiles from a single CLI flag, with no cluster, recorder, or operator prerequisites.
    \item Empirical evidence that scan-then-enforce improves post-exploitation containment relative to default \texttt{runc} on a reproducible laboratory workload.
    \item Evidence-based guidance for organisations choosing between runtime-side automation, manual policy authoring, cluster-side recorders, and full sandboxing (the last analysed from the literature).
    \item A reproducible security laboratory (\texttt{runc-hardened-test}) that documents how the empirical claims can be replayed on a stock Linux host with cgroup~v2 and AppArmor.
\end{enumerate}

\section{Thesis Organization}

Chapter~\ref{ch:litreview} reviews the relevant literature and the competitive landscape from manual authoring to full sandboxing. Chapter~\ref{ch:design} describes the design goals, the threat model, and the scanner architecture. Chapter~\ref{ch:impl} covers the implementation. Chapter~\ref{ch:experiments} reports the experimental setup and results. Chapter~\ref{ch:discussion} interprets those results and states threats to validity. Chapter~\ref{ch:conclusion} concludes and outlines future work.
